% ===================================================================
%  तक्ता क्रमांक १४ — रस्ता. --- व्यापारी.
%
%  Traduction, faite en 2026, en marathi d'aujourd'hui, sur l'ido de
%  texto/io. Regles de la colonne : voir 10-takta-01.tex. Une seule
%  suite d'alineas numerotes : les cles ne portent pas d'indice de
%  scene.
%
%  LE PLUS LONG DEFILE DE METIERS DU LIVRET, ET LE MARATHI EN NOMME
%  LA PLUPART PAR LEUR CASTE, comme aux tableaux 5, 7 et 13 :
%
%    सोनार        l'orfevre        सराफ (29)    le changeur
%    कल्हईवाला (14) le ferblantier   न्हावी (23)  le coiffeur
%    मिठाईवाला (31) le patissier     फेरीवाला (69) le colporteur
%    फुलवाली (92)  la bouquetiere   परटीण        la blanchisseuse
%    दाई (98)     la bonne d'enfants
%
%  ET LA MUSIQUE MILITAIRE EST MARATHE JUSQU'A L'INSTRUMENT. Le
%  clairon (44) est तुतारी — la trompe des armees de Shivaji —, le
%  regiment (41) est पलटण, et l'harmonium du magasin de musique n'a
%  meme pas a etre traduit : c'est l'instrument central de la
%  musique marathe depuis le XIXe siecle, arrive de France et reste
%  ici quand l'Europe l'a laisse tomber. Un emprunt qui a mieux
%  vieilli chez l'emprunteur que chez le preteur.
%
%  ET LA LYRE DE L'ENSEIGNE EST वीणा, qui n'est pas une lyre : c'est
%  l'instrument d'ici, et c'est ainsi qu'un enseigne dorée en forme
%  de lyre se lit au Maharashtra. La chose reste une enseigne de
%  marchand de musique ; c'est le NOM qui change de pays, et c'est
%  la seule chose qu'on ait le droit de faire changer.
%
%  DEUX MOTS ONT DU ETRE SEPARES QUE LE MARATHI CONFOND. वकील est a
%  la fois le consul et l'avocat, et les deux sont dans ce tableau
%  (9 et 82) : on garde वकिलात pour le consulat — le mot est vieux
%  et juste — et l'on ecrit कॉन्सुल pour l'homme, वकील pour
%  l'avocat. De meme रंगारी, le peintre du tableau 5, ne peut pas
%  servir au teinturier (53).
%
%  TROIS INVERSIONS PREVENUES — le bureau du cousin, le jet d'eau de
%  la fontaine, l'imperiale de l'omnibus — ET TROIS SEULEMENT QUI
%  ONT PASSE, sur quarante-deux blocs et pres de cent renvois. C'est
%  le meilleur rapport de toute la colonne, et il confirme ce que le
%  tableau 13 avait montre a l'envers : ce qui coute, ce n'est pas
%  la longueur d'un tableau mais sa syntaxe. Celui-ci enumere des
%  metiers cote a cote ; le 13 accrochait des choses les unes aux
%  autres, et coutait neuf inversions pour un tiers du texte.
% ===================================================================

%%K t14-apar-1 apar
\VUsaut{4.0mm}
\VUtitre{15.0pt}{60}{तक्ता क्रमांक १४}
\VUsaut{3.0mm}

%%K t14-tit-1 sub
\VUpk{11.4pt}{40}{रस्ता. --- व्यापारी.}
\VUsaut{3.0mm}

%%K t14-01-1 p
१. --- गावात राहणारा माझा मित्र जॉन माझ्या घरी तीन दिवस घालवायला आला. \VUgras{तोपर्यंत} त्याला
मोठे शहर पाहायची संधी मिळाली नव्हती; म्हणून आम्ही आमचे सगळे दिवस फिरण्यात घालवतो, आणि त्याला
माहीत नसलेल्या गोष्टी मी त्याला समजावून सांगतो.

%%K t14-02-1 p
२. --- आधी आम्ही \VUgras{वेधशाळा} \textsuperscript{(1)} पाहायला गेलो, जी त्याला फार रोचक
वाटली. ज्या \VUgras{रस्त्यावर} \textsuperscript{(3)} \VUgras{तुर्की स्नानगृह}
\textsuperscript{(2)} आहे त्या रस्त्याने परतताना आम्हाला एक \VUgras{अंत्ययात्रा}
\textsuperscript{(4)} भेटली. \VUgras{अंत्ययात्रेच्या} पुढे \VUgras{धर्मगुरूंचा वर्ग} चालत
होता, \VUgras{शववाहिकेच्या} पुढे, जिच्यात \VUgras{शवपेटी} ठेवली होती. पुष्कळ मित्र
\VUgras{शोकाकुल} कुटुंबासोबत होते.

%%K t14-02-2 p
अंत्ययात्रा गेल्यावर आम्ही त्या मोठ्या \VUgras{बिअरगृहासमोरून} \textsuperscript{(5)} रस्ता
ओलांडला, ज्याच्या गच्चीवर पुष्कळ \VUgras{गिऱ्हाइके} \VUgras{बिअर} पीत होती, काचेच्या
\VUgras{छपराखाली} \textsuperscript{(6)} सावलीत.

%%K t14-03-1 p
३. --- \VUgras{मद्यविक्रेत्याच्या} \textsuperscript{(7)} दुकानात आणि
\VUgras{संगीतविक्रेत्याच्या} \textsuperscript{(10)} दुकानात लावलेल्या \VUgras{बोधचिन्हाने} जॉन
थक्क झाला. त्याचा अर्थ काय असे त्याने मला विचारले; मी सांगितले की ते इंग्रज \VUgras{वकिलात}
\textsuperscript{(8)} दाखवते, आणि सज्ज्यावर उभे असलेले \VUgras{कॉन्सुल} \textsuperscript{(9)}
त्याला दाखवले.

%%K t14-tit-2 sub
\VUpk{10.2pt}{40}{दुकानांची पाहणी.}
\VUsaut{3.0mm}

%%K t14-04-1 p
४. --- अर्थातच आम्ही \VUgras{वाद्यांच्या} मांडणीसमोर थांबलो — \VUgras{हार्मोनियम},
\VUgras{ऑर्गन} आणि \VUgras{पियानो}; \VUgras{गाण्यांच्या वह्यांची} आणि पियानोसाठीच्या
\VUgras{स्वरलिपींची} चित्रे असलेली मुखपृष्ठे आम्ही पाहिली, आणि पाटी म्हणून वापरलेली सोन्याचा
मुलामा दिलेली लाकडी \VUgras{वीणा} न्याहाळली.

%%K t14-05-1 p
५. --- \VUgras{झाडूवाल्यांनी} \textsuperscript{(11)} आणि \VUgras{झाडणाऱ्या गाडीने}
\textsuperscript{(12)} उडवलेल्या धुळीच्या ढगांमुळे आम्हाला \VUgras{फर्निचरविक्रेत्याच्या}
\textsuperscript{(13)} सुंदर \VUgras{फर्निचरसमोरून} आणि \VUgras{कल्हईवाल्याच्या}
\textsuperscript{(14)} \VUgras{शेगड्यांच्या} व \VUgras{दिव्यांच्या} मांडणीसमोरून भराभर जावे
लागले.

%%K t14-06-1 p
६. --- शिवाय, इथे आम्हाला पदपथ सोडावा लागला, कारण तो गाठोड्यांनी अडलेला होता, जी
\VUgras{मालवाहू गाडीतून} \textsuperscript{(16)} काढून नुकत्याच भाड्याने घेतलेल्या एका
\VUgras{दुकानात} \textsuperscript{(15)} ठेवायची होती.

%%K t14-06-2 p
त्यामुळे आम्हाला \VUgras{गाडी बनवणाऱ्याच्या} \textsuperscript{(17)} सुंदर गाड्या पाहता आल्या
नाहीत, पण \VUgras{खोकी बनवणाऱ्याला} \textsuperscript{(19)} आपल्या शेजाऱ्यासाठी —
\VUgras{सायकल आणि मोटार विकणाऱ्यासाठी} \textsuperscript{(18)} — खोके तयार करताना पाहायला
थांबण्यापासून काही आम्हाला अडवले नाही. मग बराच उशीर झाल्यामुळे आम्ही घरी परतलो.

%%K t14-tit-3 sub
\VUpk{10.2pt}{40}{शहरातून फेरफटका.}
\VUsaut{3.0mm}

%%K t14-07-1 p
७. --- दुसऱ्या दिवशी माझ्या आईने जॉनला सुचवले की त्याचे छायाचित्र काढून घ्यावे, आणि मग त्याला
\VUgras{छायाचित्रकाराकडे} \textsuperscript{(20)} घेऊन जायला मला सांगितले. न्याहारीनंतर आम्ही
त्या कलावंताच्या \VUgras{कार्यशाळेत} गेलो, जिथे आम्हाला बराच वेळ थांबावे लागले. वेळ घालवायला
आम्ही भिंतीवर टांगलेल्या \VUgras{तसबिरी} \textsuperscript{(22)} पाहिल्या.

%%K t14-07-2 p
आमची पाळी आली तेव्हा काम फार वेळ चालले नाही, आणि माझ्या मित्राने \VUgras{छायाचित्रयंत्रासमोर}
\textsuperscript{(21)} बसणे संपवताच आम्ही खाली उतरलो.

%%K t14-08-1 p
८. --- पहिल्या मजल्यावर, \VUgras{न्हाव्याच्या} \textsuperscript{(23)} उघड्या दारातून आम्हाला
त्याचा नोकर एका गिऱ्हाइकाचे केस कापताना दिसला.

%%K t14-08-2 p
रस्त्यावर आल्यावर आम्ही \VUgras{तंबाखूच्या दुकानासमोरून} \textsuperscript{(24)} गेलो. लाल
\VUgras{कंदील} \textsuperscript{(25)} आणि \VUgras{पाटी} \textsuperscript{(26)} म्हणून वापरलेली
मोठी \VUgras{चिलीम} पाहून जॉनची इतकी करमणूक झाली की \VUgras{तपकीर ओढत} \textsuperscript{(27)}
गप्पा मारणाऱ्या दोन म्हाताऱ्या गृहस्थांना तो धडकला.

%%K t14-tit-4 sub
\VUpk{10.2pt}{40}{मिठाईचे दुकान.}
\VUsaut{3.0mm}

%%K t14-09-1 p
९. --- \VUgras{सराफाच्या} \textsuperscript{(29)} \VUgras{दर्शनी काचेत} मांडलेल्या सगळ्या
देशांच्या \VUgras{नोटांनी} आणि \VUgras{नाण्यांनी} काही क्षण आमचे लक्ष वेधून घेतले, पण खाऊची
वेळ झाल्यामुळे आम्ही अधिक न थांबता \VUgras{मिठाईवाल्याकडे} \textsuperscript{(31)} शिरलो.

%%K t14-10-1 p
१०. --- दुकानात असलेले सगळे \VUgras{चवदार पदार्थ} पाहून — \VUgras{केक}, \VUgras{मधाचे केक},
\VUgras{बिस्किटे} आणि सगळ्या प्रकारची \VUgras{मिठाई} — आम्हाला निवड करणे सोपे गेले नाही. शेवटी
जॉनने \VUgras{मलईचे केक} निवडले, आणि मी फक्त एक लहानशी \VUgras{चेरीची टार्ट} घेतली, कारण
\VUgras{गोड पदार्थांनी माझे दात दुखतात}, आणि \VUgras{दंतवैद्याकडे} \textsuperscript{(30)} फार
वेळा जायची मला मुळीच इच्छा नाही.

%%K t14-tit-5 sub
\VUpk{10.2pt}{40}{टंकलेखन.}
\VUsaut{3.0mm}

%%K t14-11-1 p
११. --- माझ्या मित्राला \VUgras{टंकलेखनयंत्र} दाखवायला — ज्याबद्दल त्याने ऐकले होते पण जे
त्याने पाहिले नव्हते — आम्ही ते \VUgras{कार्यालय} \textsuperscript{(32)} गाठले, जे माझ्या
\VUgras{चुलतभावाचे} \textsuperscript{(34)} आहे. तो आपली पत्रे आपल्या \VUgras{सचिवाला}
\textsuperscript{(35)} सांगत होता, जो \VUgras{लघुलेखक} आहे. पत्र संपले न संपले तोच
\VUgras{टंकलेखकाने} \textsuperscript{(33)} ते आपल्या यंत्रावर उतरवले. माझ्या नातेवाइकाच्या
कामात व्यत्यय नको म्हणून आम्ही त्याच्याकडे थोडाच वेळ थांबलो.

%%K t14-12-1 p
१२. --- माझ्या चुलतभावाच्या कार्यालयाखाली एक मोठा \VUgras{घड्याळे आणि दागिने विकणारा}
\textsuperscript{(37)} राहतो, जो शिवाय \VUgras{सोनारही} आहे. त्याच्या भव्य दुकानात पुष्कळ
\VUgras{घड्याळे}, \VUgras{लंबकाची घड्याळे}, \VUgras{खिशातली घड्याळे}, \VUgras{साखळ्या} आणि
मौल्यवान \VUgras{दागिने} आहेत — उदाहरणार्थ \VUgras{मोत्यांचे हार}, सोन्याची \VUgras{कडी},
\VUgras{कर्णफुले} आणि \VUgras{रत्नांनी} व \VUgras{हिऱ्यांनी} जडवलेल्या \VUgras{अंगठ्या}. एक
दर्शनी काच खास \VUgras{सोन्याच्या कामासाठी} राखीव आहे आणि तिच्यात
\VUgras{चांदीच्या भांड्यांचा} मोठा संग्रह आहे.

%%K t14-13-1 p
१३. --- रस्त्याच्या कोपऱ्यावर वळताना माझ्या लक्षात आले की घराच्या \VUgras{क्रमांकाजवळ} लावलेली
\VUgras{पाटी} \textsuperscript{(36)} बदलली आहे. मी माझे निरीक्षण मोठ्याने सांगणार होतो,
तेवढ्यात \VUgras{त्या} लष्करी वाद्यवृंदाचे आनंदी सूर ऐकू आले. आम्ही पलटणीच्या दिशेने धावत
सुटलो, हे न विचारता की \VUgras{ती} \VUgras{टोपीविक्रेत्याच्या} \textsuperscript{(39)} आणि
\VUgras{हातमोजेविक्रेत्याच्या} \textsuperscript{(40)} दुकानांसमोरून जाणार आहे, आणि तिचे संचलन
पाहायला आम्ही \VUgras{चष्मेवाल्याच्या} \textsuperscript{(38)} दर्शनी काचेसमोर थांबूनही तेवढेच
चांगले उभे राहिलो असतो — ती काच \VUgras{नाकावरच्या चष्म्यांनी},
\VUgras{कानाला अडकवायच्या चष्म्यांनी} आणि \VUgras{दुर्बिणींनी} पुरी भरलेली आहे.

%%K t14-tit-6 sub
\VUpk{10.2pt}{40}{संचलन.}
\VUsaut{3.0mm}

%%K t14-14-1 p
१४. --- \VUgras{पलटणीच्या} \textsuperscript{(41)} पुढे \VUgras{ढोलप्रमुख}
\textsuperscript{(42)} चालत होता, त्याच्यामागे \VUgras{ढोलकरी} \textsuperscript{(43)},
\VUgras{तुतारीवाले} \textsuperscript{(44)} आणि सगळा \VUgras{वादकवृंद}. आपल्या सहायकासोबत चौकात
असलेला \VUgras{सेनापती} \textsuperscript{(45)} \VUgras{संचलन} पाहायला \VUgras{तोटीजवळ}
\textsuperscript{(49)} थांबला होता — \VUgras{कारंज्याच्या} \textsuperscript{(48)}.

%%K t14-14-2 p
\VUgras{मुख्यालयाच्या अधिकाऱ्यांनी} \textsuperscript{(46)}, \VUgras{कर्नलने} आणि \VUgras{उपकर्नलने} त्याला तलवारीने मुजरा केला. \VUgras{मेजर} आपापल्या \VUgras{पलटणींपुढे} \textsuperscript{(47)} होते आणि प्रत्येक \VUgras{कप्तान} आपल्या \VUgras{तुकडीला} हुकूम देत होता, तर \VUgras{लेफ्टनंट}, \VUgras{उपलेफ्टनंट} आणि \VUgras{उपअधिकारी} आपापल्या माणसांशेजारी नेहमीच्या जागी होते.

%%K t14-15-1 p
१५. --- पुष्कळ वाटसरू \VUgras{चौरस्त्यावर} \textsuperscript{(50)} गोळा झाले होते; व्यापारी,
\VUgras{छत्रीवाला} \textsuperscript{(51)}, \VUgras{कपडे रंगवणारा} \textsuperscript{(53)},
\VUgras{शस्त्रविक्रेता} \textsuperscript{(54)} \VUgras{शिपाई} पाहायला बाहेर आले. सगळी वाहने —
\VUgras{टांगे} \textsuperscript{(52)}, \VUgras{सामानाच्या गाड्या} \textsuperscript{(57)} आणि
\VUgras{पाणी शिंपडणारी गाडीसुद्धा} \textsuperscript{(56)} — पदपथाच्या कडेला उभी राहिली.

%%K t14-tit-7 sub
\VUpk{10.2pt}{40}{रस्त्यावरचे प्रसंग.}
\VUsaut{3.0mm}

%%K t14-15-2 p
हातात आपल्या \VUgras{नमुन्यांच्या पेट्या} घेतलेला एक \VUgras{व्यापारी प्रवासी}
\textsuperscript{(55)} भराभर रस्ता ओलांडून गेला, आणि \VUgras{वरच्या मजल्यावर}
\textsuperscript{(59)} बसलेली माणसे — \VUgras{घोडागाडीच्या} \textsuperscript{(58)} — तो देखावा
पाहायला वळली. एका बिचाऱ्या \VUgras{भटक्या गवयाला} \textsuperscript{(60)}, जो आपला
\VUgras{फिरत्या पेटीचा वाजा} \textsuperscript{(61)} घेऊन आला होता, त्याला आपले \VUgras{गाणे}
थांबवावे लागले, कारण ढोलांच्या गजरात त्याचा आवाज बुडून गेला.

%%K t14-16-1 p
१६. --- पलटण \VUgras{कापडाच्या दुकानासमोर} \textsuperscript{(64)} आली तेव्हा
\VUgras{शिवणारणींनी} \textsuperscript{(63)} आणि \VUgras{कापणाऱ्यांनी} \textsuperscript{(62)}
आपली \VUgras{शिवणयंत्रे} आणि आपले काम सोडून खिडकीतून पाहिले. \VUgras{व्यापाऱ्याला}
\textsuperscript{(65)}, ज्याने नुकतेच नवे \VUgras{कपडे} \textsuperscript{(66)} आपल्या
\VUgras{मांडणीत} लावले होते, त्याला \VUgras{कापडे} \textsuperscript{(67)} आणि \VUgras{तागे} आत
ठेवावे लागले — पदपथावर \VUgras{गटाराच्या झाकणाजवळ} \textsuperscript{(68)} मांडलेले —, गर्दी ती
पाडील या भीतीने; एका म्हाताऱ्या \VUgras{फेरीवाल्यालाही} \textsuperscript{(69)} — ज्याच्याकडून
एक द्वारपालीण मोजांचा भाव करत होती — आपली विक्री पुरी करायला एका बोळात शिरावे लागले.

%%K t14-16-2 p
आम्ही पलटणीसोबत छावणीपर्यंत गेलो, \VUgras{आणि}, आमच्या दिवसाबद्दल फार समाधानी होऊन जेवणाच्या
वेळेला घरी परतलो.

%%K t14-tit-8 sub
\VUpk{10.2pt}{40}{निरनिराळे व्यापार आणि व्यवसाय.}
\VUsaut{3.0mm}

%%K t14-17-1 p
१७. --- दुसऱ्या दिवशी माझ्या वडिलांनी आम्हाला तिथल्या वर्तमानपत्राचा छापखाना पाहायला सुचवले.
आम्ही उत्साहाने होकार दिला.

%%K t14-17-2 p
\VUgras{छापणारा} हा \VUgras{पुस्तकविक्रेताही} \textsuperscript{(70)} \VUgras{(म्हणजे पुस्तके विकणारा),} \VUgras{प्रकाशक}, \VUgras{कागदविक्रेता} आणि \VUgras{पुस्तक बांधणाराही} आहे. त्याने पहिल्या मजल्यावर वर्तमानपत्राचे \VUgras{संपादकीय कार्यालय} \textsuperscript{(71)} थाटले आहे, आणि \VUgras{कोरीवकामाची खोलीही}, जिथे \VUgras{शिलामुद्रक} \textsuperscript{(72)} आणि \VUgras{तांब्यावर कोरणारे} \textsuperscript{(73)} काम करतात, आणि शेवटी \VUgras{जुळणीची खोली}, जिथे \VUgras{खिळेजुळणी करणारे} \textsuperscript{(74)} असतात. खरी \VUgras{छपाईची खोली} \textsuperscript{(75)}, \VUgras{छापयंत्रे} आणि निरनिराळी यंत्रे तळमजल्यावर आहेत.

%%K t14-tit-9 sub
\VUpk{10.2pt}{40}{जॉन खूप प्रश्न विचारतो.}
\VUsaut{3.0mm}

%%K t14-18-1 p
१८. --- छापखान्यातून बाहेर पडताना जॉन फार थक्क होऊन \VUgras{सुतासुयांच्या दुकानासमोरच्या}
\textsuperscript{(77)} एका खांबावर बसवलेल्या काचेच्या लहानशा पेटीसमोर थांबला आणि ती कशासाठी
वापरतात असे विचारले. माझ्या वडिलांनी त्याला समजावले की ती \VUgras{आगीची इशारा-पेटी}
\textsuperscript{(76)} आहे. एरवीही शहरातले सगळेच माझ्या मित्राला थक्क करणारे होते — साध्या
\VUgras{दगडी कारंज्यापासून} \textsuperscript{(78)} आणि हलक्या \VUgras{मालवाहू तीनचाकीपासून}
\textsuperscript{(80)}, जी गणवेश घातलेला एक \VUgras{पोऱ्या} \textsuperscript{(79)} चालवत होता,
\VUgras{टपालगाडीपर्यंत} \textsuperscript{(81)}.

%%K t14-19-1 p
१९. --- माझे वडील काही क्षण आम्हाला सोडून \VUgras{करवसुली अधिकाऱ्याशी} \textsuperscript{(83)}
बोलायला गेले — जो एका \VUgras{वकिलाशी} \textsuperscript{(82)} आणि एका \VUgras{नोटरीशी}
\textsuperscript{(84)} गप्पा मारायला थांबला होता — तेवढ्यात आम्ही रस्त्यावरची वर्दळ डोळ्यांनी
न्याहाळत करमणूक करून घेतली. अत्यंत निरनिराळी माणसे आमच्यासमोरून गेली : टोपलीत भव्य
\VUgras{पेस्ट्री} घेऊन जाणारा एक \VUgras{मिठाईवाल्याचा नोकर} \textsuperscript{(85)} एका
\VUgras{कपडेविक्रेत्यामागून} \textsuperscript{(86)} आला; एक \VUgras{दिवे लावणारा}
\textsuperscript{(87)} एका \VUgras{गॅसवाल्याशी} \textsuperscript{(88)} बोलत होता; एक
\VUgras{धर्मभगिनी} \textsuperscript{(89)} अनाथ मुलीचा हात धरून आपल्या मठाकडे परतत होती.

%%K t14-tit-10 sub
\VUpk{10.2pt}{40}{घरी परतताना.}
\VUsaut{3.0mm}

%%K t14-20-1 p
२०. --- माझे वडील आम्हाला येऊन मिळाले त्या क्षणी, काजळीने पुरत्या माखलेल्या दोन
\VUgras{धुराडी साफ करणाऱ्यांचे} \textsuperscript{(91)} मळके चेहरे आणि विचित्र फाटके कपडे पाहून
जॉन खो खो हसला.

%%K t14-21-1 p
२१. --- मला माझ्या आईसाठी \VUgras{फुलवालीकडून} \textsuperscript{(92)} एक रोप घ्यायचे होते, पण
माझ्या वडिलांनी मला दाखवून दिले की ते घेऊन जायला फार जड होईल आणि \VUgras{संत्री} घेणे बरे.
आम्ही \VUgras{विक्रेतीणीजवळ} \textsuperscript{(94)} गेलो, आणि आपल्या मुलासाठी निवड करत
असलेल्या \VUgras{संगोपिकेची} \textsuperscript{(93)} खरेदी संपल्यावर आम्ही आमचे घेतले.

%%K t14-22-1 p
२२. --- घरी परतताना आम्हाला कित्येक ओळखीची माणसे भेटली : माझ्या घरची एक म्हातारी मैत्रीण, जिने
आपल्या \VUgras{सोबतिणीच्या} \textsuperscript{(95)} दंडाचा आधार घेतला होता, पूल आणि रस्त्यांचे
\VUgras{अभियंता} \textsuperscript{(96)}, आणि माझी एक चुलतबहीण आपल्या \VUgras{दाईसह}
\textsuperscript{(98)}.

%%K t14-22-2 p
मग आम्हाला माझ्या आईची \VUgras{टोप्या बनवणारी} \textsuperscript{(97)} भेटली, आणि शहरात परक्या
असलेले दोन \VUgras{भिक्षु} \textsuperscript{(99)}, जे स्थानकाकडे जाणारी वाट विचारायला माझ्या
वडिलांकडे आले.

\VUsaut{6.0mm}
\VUfilet{22mm}
